\documentclass[11pt,a4paper]{article}

\usepackage[margin=25mm]{geometry}
\usepackage{amsmath,amssymb,mathtools}
\usepackage{booktabs}
\usepackage{enumitem}
\usepackage[hidelinks]{hyperref}

\newcommand{\R}{\mathbb{R}}
\newcommand{\N}{\mathbb{N}}
\newcommand{\Oh}{\mathcal{O}}
\setlength{\emergencystretch}{2em}

\title{EE6407: Genetic Algorithms and Machine Learning\\
       \large Quiz 1 Review Notes --- Evolutionary Computing Part 1
       (Weeks 1--4)}
\author{Living revision notes}
\date{Updated 24 August 2026}

\begin{document}
\maketitle

\paragraph{Sources and scope.}
This note covers the complete evolutionary-computing Part~1 material in
\path{../resources/lecture-01-evolutionary-computing-part-1.pdf}
(36 PDF pages, 72 numbered slides), supported by the Week~1 and Week~2 class
transcripts in \path{../resources/week-01-transcript.txt} and
\path{../resources/week-02-transcript.txt}. The specialised $N$-queens method
is based on Gwee and
Lim, ``An evolution search algorithm for solving N-queen problems,'' supplied
as \path{../resources/n-queens-evolution-search-paper.pdf} (especially
article pp.~1--4). The latter is a
particular evolutionary search algorithm, not the only valid genetic
algorithm for the problem.

\paragraph{Assessment facts visible in the supplied lecture slides.}
Continuous assessment contributes 40\% and the three-hour final examination
contributes 60\%.  For the first four-week part, one closed-book written quiz
contributes 10\% and is scheduled for the last hour of Week 4.  Exact date,
time, room allocation, and any later changes should be checked in NTULearn.

\tableofcontents

\section{Why Evolutionary Computing?}

Evolutionary computing (EC) turns a simplified metaphor of biological
evolution into a family of computational search methods.  A population of
candidate solutions competes under a task-specific fitness measure.  Selected
candidates create new candidates through variation, and the process repeats.
Genetic algorithms (GAs) are one historically important member of this wider
family.

The attraction is practical: many engineering, design, scheduling, and
learning problems have a huge or irregular search space, mixed discrete and
continuous decisions, non-differentiable objectives, or many local optima.
An evolutionary algorithm (EA) does not require gradients and can search with
multiple candidates at once.  This flexibility does \emph{not} imply a
guarantee of finding the global optimum.

\subsection{Historical landmarks}

The lecture highlights Turing's proposal of genetical or evolutionary search
(1948), Bremermann's optimisation by evolution and recombination (1962),
Rechenberg's evolution strategies (1964), Fogel, Owens, and Walsh's
evolutionary programming (1965), Holland's genetic algorithms (1975), and
Koza's genetic programming (1992).  These traditions originally emphasised
different representations, but modern practice starts from the problem and
chooses a suitable representation and operators.

\section{Classifying Problems Before Solving Them}

The purpose of classification is not taxonomy for its own sake.  It exposes
what is unknown, how large the search is, what counts as a good answer, and
what computational difficulty should be expected.

\subsection{The input--model--output view}

Represent a system abstractly by
\[
  y=M(x),
\]
where $x$ is the input, $M$ the model, and $y$ the output.  Holding two parts
fixed and seeking the third gives three different problem types.

\begin{center}
\begin{tabular}{lll p{0.38\textwidth}}
\toprule
Unknown & Given & Problem type & Typical question \\
\midrule
$x$ & $M,y$ & optimisation & Which design/input gives the desired output? \\
$M$ & paired $x,y$ data & modelling & Which mapping explains or predicts the data? \\
$y$ & $x,M$ & simulation & What happens under this input and model? \\
\bottomrule
\end{tabular}
\end{center}

Examples include university timetabling, satellite structural design, the
travelling-salesman problem (optimisation), credit-risk prediction and voice
control (modelling), and weather or policy ``what-if'' studies (simulation).
A modelling task can often be converted into optimisation by parameterising
$M_\theta$ and minimising a prediction loss over $\theta$.

\subsection{Search spaces and problem solvers}

The \emph{search space} $S$ is the set of all candidate objects under
consideration; the feasible set $\mathcal{F}\subseteq S$ contains candidates
that satisfy every hard constraint.  A problem definition describes $S$,
$\mathcal{F}$, and the desired quality.  A problem solver describes how to
move through or sample $S$.  Keeping these notions separate prevents a common
mistake: a large search space is a property of the formulation, while slow
progress may also result from an unsuitable algorithm or representation.

For a tour through $n$ labelled cities, the naive number of orderings is
$n!$ (or $(n-1)!/2$ if rotations and reversal are identified for a symmetric
tour).  Exhaustive enumeration therefore becomes impractical very quickly.

\subsection{Optimisation and constraint satisfaction}

An objective function $f:S\to\R$ ranks candidate quality.  A constraint is a
condition such as $g_j(x)\leq 0$ or $h_k(x)=0$.

\begin{itemize}[leftmargin=*]
  \item A \emph{free optimisation problem} has an objective but no explicit
        constraints.
  \item A \emph{constraint-satisfaction problem} (CSP) asks only for a
        feasible point.
  \item A \emph{constrained optimisation problem} (COP) asks for the best
        feasible point.
\end{itemize}

A CSP can be transformed into optimisation by defining a violation measure,
for example
\[
  V(x)=\sum_j \max\{0,g_j(x)\}
       +\sum_k |h_k(x)|,
\]
and minimising $V$.  A COP can use a penalty objective
\[
  F_\lambda(x)=f(x)+\lambda V(x)
\]
for a minimisation problem.  The penalty weight $\lambda$ must be chosen with
care: too small allows attractive but infeasible candidates; too large can
flatten useful distinctions among infeasible candidates.

\subsection{P, NP, NP-complete, and NP-hard}

Let the input size be $n$, measured by the number of bits or symbols required
to encode an instance.  Running time is usually discussed in the worst case
as a function of $n$.

\begin{itemize}[leftmargin=*]
  \item $\mathbf{P}$: decision problems solvable in polynomial time.
  \item $\mathbf{NP}$: decision problems whose proposed yes-certificates can
        be verified in polynomial time.  Thus $\mathbf{P}\subseteq\mathbf{NP}$.
  \item $\mathbf{NP}$-complete: problems in $\mathbf{NP}$ to which every
        problem in $\mathbf{NP}$ can be reduced in polynomial time.
  \item $\mathbf{NP}$-hard: problems at least as hard as the
        $\mathbf{NP}$-complete problems; they need not themselves be decision
        problems or belong to $\mathbf{NP}$.
\end{itemize}

It remains unknown whether $\mathbf{P}=\mathbf{NP}$.  Metaheuristics such as
EAs are therefore used to obtain good solutions within a practical budget,
especially when exact solution is too expensive.  Their use does not change
the complexity class of the underlying problem.

\section{From Natural Evolution to an Algorithm}

The simplified Darwinian metaphor has four essential ideas:

\begin{enumerate}[leftmargin=*]
  \item resources are finite, so selection is unavoidable;
  \item individuals differ in phenotype;
  \item some differences affect reproductive success and are heritable;
  \item random variation continually introduces diversity.
\end{enumerate}

Individuals are the units on which selection acts, but the population is the
unit that evolves.  Computationally, fitness is an explicitly designed score,
unlike natural fitness, which is inferred indirectly from reproductive
success.

\subsection{Genotype, phenotype, genes, and alleles}

The \emph{phenotype} is a solution in the original problem domain.  The
\emph{genotype} is the manipulable encoding used by the EA.  A chromosome is
composed of genes at loci, and each gene takes an allele value.

Let $e$ be an encoder and $d$ a decoder:
\[
  e:\mathcal{F}\to G, \qquad d:G\to S.
\]
The decoder must give an unambiguous candidate for every admissible genotype.
For global search, every feasible phenotype should be reachable from at least
one genotype.  Encoding need not be one-to-one, but redundancy can bias the
search because phenotypes with many encodings are sampled more often.

\paragraph{Representation principle.}
Operators must preserve or sensibly repair the chosen representation.  Bit
flip suits binary strings; Gaussian perturbation suits real vectors; swap or
insertion mutations suit permutations.  Using ordinary one-point crossover
on a permutation may duplicate some values and omit others, producing an
invalid solution.

\section{The General Evolutionary Algorithm}

An EA is a stochastic, population-based, generate-and-test algorithm.
One generic form is:

\begin{enumerate}[leftmargin=*]
  \item initialise a population $P_0$ of candidate genotypes;
  \item decode and evaluate every candidate;
  \item while the termination condition is false:
    \begin{enumerate}
      \item select parents from $P_t$;
      \item recombine selected parents;
      \item mutate the resulting offspring;
      \item evaluate the offspring;
      \item select survivors to form $P_{t+1}$.
    \end{enumerate}
\end{enumerate}

The process balances two opposing forces:
\[
  \underbrace{\text{mutation + recombination}}_{\text{novelty/diversity}}
  \qquad\text{against}\qquad
  \underbrace{\text{parent + survivor selection}}_{\text{quality pressure}}.
\]
Excessive selection pressure can cause premature convergence; insufficient
pressure wastes evaluations on poor regions.  Diversity is therefore not an
end in itself but a resource that keeps alternatives available.

\subsection{Fitness and population}

The fitness or evaluation function maps each phenotype to a scalar used for
comparison.  It must reflect the actual task, including constraint handling,
and discriminate usefully among candidates.  If the natural objective is a
cost $c(x)$ to minimise, it is often safer to keep a minimisation-aware
selection method or use a monotone transformation rather than the singular
$1/c(x)$ when $c(x)$ may be zero.

A population is a multiset, so duplicate genotypes are allowed.  Diversity
can refer to genotype, phenotype, or fitness diversity; these are different.
For example, several genotypes may decode to the same phenotype, and distinct
phenotypes may have equal fitness.

\subsection{Parent and survivor selection}

Selection favours quality without normally making the outcome certain.  In
fitness-proportionate (roulette-wheel) selection, assuming non-negative
fitness,
\[
  \Pr(i\text{ selected})=\frac{f_i}{\sum_{j=1}^{\mu} f_j}.
\]
Thus candidates $A,B,C$ with fitness $3,1,2$ receive probabilities
$1/2,1/6,1/3$.  Scaling matters: adding a constant to all fitness values
changes these probabilities even though the ranking is unchanged.  Rank or
tournament selection avoids this particular sensitivity.

Survivor selection controls replacement.  Common choices include deleting
all parents after producing $\mu$ offspring, keeping the best $\mu$ among
parents and offspring, or using \emph{elitism} to guarantee that a small
number of the best candidates survive.  Strong elitism preserves progress but
can reduce diversity.

\subsection{Variation}

Mutation is unary and makes a random change to one genotype.  It introduces
or preserves diversity and, when the mutation graph connects the whole
genotype space, makes every solution reachable in principle.

Recombination has arity greater than one, usually two, and merges information
from parents.  A child is not guaranteed to be better than either parent; the
hope is that useful partial structures can occasionally be combined.

\subsection{Initialisation and termination}

Random initialisation should cover the admissible allele values without
unintended bias.  Known solutions or domain heuristics may seed the
population, but excessive seeding can suppress exploration.

Typical stopping rules are a maximum generation or evaluation count, reaching
a target fitness, no meaningful improvement for several generations, or
diversity falling below a threshold.  Stopping because the population is
unchanged is evidence of stagnation, not proof of global optimality.

\section{Worked Model: The \texorpdfstring{$N$}{N}-Queens Problem}

Place $N$ queens on an $N\times N$ board so that no two share a row, column,
or diagonal.  With exactly one queen assigned to each row, encode a candidate
as a permutation
\[
  q=(q_1,q_2,\ldots,q_N),
\]
where $q_i\in\{1,\ldots,N\}$ is the column of the queen in row $i$.  The
permutation property gives
\[
  q_i\neq q_j \quad (i\neq j),
\]
so row and column conflicts are impossible by construction.  Only diagonal
conflicts remain:
\[
  |q_i-q_j|\neq |i-j| \quad\text{for every }i\neq j.
\]
This illustrates a major design lesson: a good representation can remove hard
constraints from the search rather than repeatedly penalising them.

\subsection{Fitness formulations}

Define a conflict indicator for a pair $i<j$:
\[
  a_{ij}(q)=
  \begin{cases}
    1, & |q_i-q_j|=|i-j|,\\
    0, & \text{otherwise}.
  \end{cases}
\]
The total number of attacking pairs is
\[
  C(q)=\sum_{1\leq i<j\leq N} a_{ij}(q).
\]
We may minimise $C(q)$, or maximise
\[
  f_{\mathrm{pairs}}(q)=\binom{N}{2}-C(q).
\]
The optimum satisfies $C(q)=0$.

The supplied paper instead gives each queen a binary gene fitness
\[
  c_i(q)=
  \begin{cases}
    1, & \text{queen }i\text{ conflicts with no other queen},\\
    0, & \text{otherwise},
  \end{cases}
  \qquad F(q)=\sum_{i=1}^{N}c_i(q).
\]
Here $F(q)=N$ exactly at a solution.  This differs from counting non-attacking
pairs: both identify the same optima, but they shape the intermediate fitness
landscape differently.

\subsection{Representation-aware operators}

A valid mutation swaps the alleles at two randomly selected loci.  The result
remains a permutation.  A valid crossover can copy a prefix from one parent,
then scan the other parent cyclically and append values not already present.
This is an order-preserving construction: one parent supplies explicit
values, while the other supplies relative order.

One lecture example chooses five candidates at random and uses the best two as
parents.  A child may replace the first lower-fitness individual in a
population sorted by decreasing fitness.  These are design choices, not fixed
requirements of all GAs.

\subsection{The specialised two-stage evolutionary search}

Gwee and Lim reduce the unrestricted $N^N$ configurations to the $N!$
permutation space.  Their article reports that solution density still falls
sharply: for $N=7$ there are about $7936$ solutions per million permutations,
whereas for $N=12$ there are about $29$ per million (article p.~1).

Their algorithm has two stages:

\begin{enumerate}[leftmargin=*]
  \item \textbf{Selection and reproduction.} Evaluate and sort the
        population; retain the fitter half and create related offspring to
        replace the weaker half.  A shift-by-one reproduction maps, for
        example, $(1,3,5,2,4)$ to $(2,4,1,3,5)$.
  \item \textbf{Selection and mutation.} Identify \emph{weak genes}, namely
        queens involved in conflicts, and preferentially mutate them.  Stop
        when $F(q)=N$ or the generation limit is reached.
\end{enumerate}

Mutating only conflicting queens can oscillate between local peaks.  The
paper's remedy is to label one or more otherwise fit genes as
\emph{pseudo-weak}; mutation may then disturb a larger structure and escape
the cycle.  This is a concrete example of injecting exploration when a
domain-specific operator has become too restrictive.  The reported simulations
found instances up to $N=2000$ on a personal computer and, for the tested
settings, first solutions within 30 generations (article p.~4).  These are
empirical results for that implementation, not universal complexity bounds.

\section{Worked Model: Simple GA for \texorpdfstring{$f(x)=x^2$}{f(x)=x squared}}

Suppose $x$ is an integer in $[0,31]$ and the goal is to maximise $x^2$.
Five bits encode every feasible value:
\[
  \begin{aligned}
    g &= b_4b_3b_2b_1b_0,\\
    x=d(g) &= \sum_{k=0}^{4}2^k b_k,
    & f(g) &= d(g)^2.
  \end{aligned}
\]

A hand simulation proceeds as follows:

\begin{enumerate}[leftmargin=*]
  \item sample a population of five-bit strings;
  \item decode each string and compute $x^2$;
  \item select a mating pool using fitness proportions;
  \item pair parents and sample a one-point crossover location from the four
        gaps between bits;
  \item independently flip each bit with mutation probability $p_m$;
  \item evaluate the new population and repeat.
\end{enumerate}

For population size $4$, the population contains $20$ bits.  With
$p_m=0.001$, the expected number of mutated bits per generation is only
$20(0.001)=0.02$.  This does \emph{not} mean exactly $0.02$ bits mutate; the
actual count is random and is usually zero in a particular generation.  The
probability of no mutation is
\[
  (1-p_m)^{20}=0.999^{20}\approx 0.9802.
\]
This distinction between expectation and realised events is important when
interpreting stochastic algorithms.

\section{EAs in Optimisation Context}

For a fixed feasible set $S$, global optimisation seeks
\[
  x^*\in\arg\max_{x\in S} f(x)
  \quad\text{or}\quad
  x^*\in\arg\min_{x\in S} f(x).
\]
Exact deterministic methods may guarantee $x^*$ but require
super-polynomial time in difficult cases.  Heuristics decide which candidate
to test next and often find useful answers quickly, but ordinarily provide no
certificate of global optimality.

Hill climbing uses a neighbourhood and often terminates at a local optimum.
EAs differ through a population, stochastic selection, and multiple variation
operators, especially recombination.  Their effective neighbourhood is
operator-dependent: mutation connects nearby genotypes, while crossover can
jump between regions defined by parental material.

\subsection{Domain knowledge and no free lunch}

Adding domain knowledge through representations, repair rules, seeded
solutions, or specialised operators can improve performance on the intended
problem class.  It simultaneously introduces bias and may hurt performance on
other classes.  The lecture's ``no free lunch'' message is therefore not that
all algorithms are equally useful on a chosen engineering problem.  It is
that, averaged uniformly over all possible objective functions under the
theorem's assumptions, gains on some problems are offset by losses on others.
The practical conclusion is to match algorithm bias to the problem and test
empirically on representative instances.

\section{Quiz 1 Rapid Review}

\paragraph{Scope boundary.}
Prepare all of Part~1 in the Lecture~1 Part~1 deck, PDF pages 1--36:
classification, complexity, EC foundations, the EA loop and components,
8-queens, the simple GA, EA behaviour, and optimisation context. Historical
questions support the patterns of classification, representation, encoding,
fitness, and hand calculations only. Differential evolution, PSO, advanced
crossover families, schema theory, niching, multi-objective optimisation, and
machine learning are not part of this Quiz~1 pack.

\subsection{Definitions and formulas to reproduce}
{\small
\begin{align*}
  x\text{ unknown}&\Rightarrow\text{optimisation},&
  M\text{ unknown}&\Rightarrow\text{modelling},\\
  y\text{ unknown}&\Rightarrow\text{simulation},&
  p_i&=\frac{f_i}{\sum_j f_j},\\
  |q_i-q_j|&=|i-j|,&
  f(q)&=\binom N2-C(q),\\
  z&=\sum_{k=0}^{L-1}2^k b_k,&
  x&=a+\frac{z}{2^L-1}(b-a),\\
  \mathbb E[\text{flips}]&=NLp_m,&
  \Pr(\text{no flips})&=(1-p_m)^{NL}.
\end{align*}
}
\noindent
$\mathbf P$ means polynomial-time solvable; $\mathbf{NP}$ means
polynomial-time verifiable yes-certificates; NP-complete means both in NP and
NP-hard; NP-hard does not require membership in NP or even a decision-problem
form. A phenotype is the domain solution, while a genotype is the encoding
manipulated by the EA.

\subsection{Algorithm steps and calculation checks}
\begin{align*}
  &\text{initialise}\to\text{evaluate}\to\text{select parents}
  \to\text{recombine}\to\text{mutate}\\
  &\hspace{28mm}\to\text{evaluate offspring}
  \to\text{select survivors}\to\text{stop/repeat}.
\end{align*}
\begin{enumerate}[leftmargin=*,nosep]
  \item Write the objective direction and larger-is-better fitness convention
        before calculating selection probabilities.
  \item Decode every chromosome after crossover and mutation; recompute
        fitness rather than carrying an old value forward.
  \item Check roulette probabilities against the correct population and
        verify that they sum to $1$ up to rounding.
  \item For permutation encodings, check for duplicates and omissions after
        every operator. Swap mutation preserves validity; ordinary one-point
        crossover generally does not.
  \item For $N$-queens, inspect each pair once with $i<j$ and use
        $\binom N2$ as the maximum fitness.
  \item Distinguish an expected mutation count from the realised random
        number of flips.
\end{enumerate}

\subsection{Likely traps}
NP does not mean non-polynomial; selection favours but does not guarantee;
parent selection and survivor selection answer different questions; high
fitness pressure can destroy diversity; convergence is not a proof of global
optimality; a local optimum depends on the chosen neighbourhood; and no-free-
lunch does not say that algorithm choice is irrelevant on a particular
problem.

\subsection{Final recall checklist}
Without looking, can you:
\begin{itemize}[leftmargin=*,nosep]
  \item classify all three $y=M(x)$ cases and identify $S$, $\mathcal F$, and
        $f$;
  \item draw the complete EA loop and state the role of each component;
  \item state the P/NP/NP-complete/NP-hard relationships without overclaiming;
  \item define genotype, phenotype, chromosome, gene, and allele;
  \item score a queen permutation and explain operator validity;
  \item decode a binary string, calculate roulette probabilities, perform one
        crossover and mutation, and verify the resulting fitnesses; and
  \item explain exploration, exploitation, local versus global optimality,
        domain bias, and the practical no-free-lunch message?
\end{itemize}

\section{Common Pitfalls}

\begin{enumerate}[leftmargin=*]
  \item \textbf{Treating GA and EC as exact synonyms.}  GA is a member of the
        broader EC family, even if the lecture sometimes uses the terms
        informally and interchangeably.
  \item \textbf{Confusing phenotype with genotype.}  Fitness belongs to the
        decoded problem solution, even when it is computed directly from a
        convenient encoding.
  \item \textbf{Using invalid operators.}  Operators designed for bit strings
        need not preserve permutations or constraints.
  \item \textbf{Assuming fitness always improves.}  Variation can produce
        worse children; improvement arises statistically through selection.
  \item \textbf{Equating convergence with proof.}  A stagnant population may
        be trapped locally.
  \item \textbf{Quoting empirical performance as a guarantee.}  Results such
        as ``under 30 generations'' depend on the paper's instances,
        parameters, operators, and implementation.
  \item \textbf{Ignoring fitness scaling.}  Roulette-wheel probabilities can
        become unstable with negative values or a single dominant value.
\end{enumerate}

\section{Revision Checklist}

You should be able to:
{\small
\begin{itemize}[leftmargin=*,nosep]
  \item classify a task using the input--model--output view;
  \item distinguish a search space, feasible set, objective, and solver;
  \item state the meanings and relationships of P, NP, NP-complete, and
        NP-hard without claiming that $\mathbf{P}\neq\mathbf{NP}$ is proved;
  \item draw the EA loop and explain the roles of representation, evaluation,
        population, selection, recombination, and mutation;
  \item explain why variation promotes novelty while selection promotes
        quality;
  \item compute roulette-wheel probabilities and discuss their sensitivity;
  \item formulate $N$-queens as a permutation and derive a conflict-based
        fitness;
  \item choose mutation and crossover operators that preserve a permutation;
  \item explain weak and pseudo-weak genes in the supplied $N$-queens paper;
  \item distinguish a good empirical result from a global-optimality or
        runtime guarantee.
\end{itemize}
}

\end{document}
